\documentclass[9pt]{extarticle}
\usepackage{core_attempt_forms}

\begin{document}

\CoreHeader{Proofs}{P01}{A--E}{\FivePartScoring Use complete sentences, name the requested method or rule, and justify every inference.}
\begin{multicols}{2}
\begin{CoreProblem}{A.4}{Inference and proof structure}
From $p\to q$, $q\to r$, and $\neg r$, derive $\neg p$. Name the rule used at
each step.
\end{CoreProblem}
\begin{CoreProblem}{B.1}{Direct proof}
Give a direct proof that if an integer $n$ is even, then $n^2$ is even.
\end{CoreProblem}
\begin{CoreProblem}{C.3}{Contraposition or contradiction}
Prove by contradiction that there is no largest integer.
\end{CoreProblem}
\columnbreak
\begin{CoreProblem}{D.1}{Cases, existence, or counterexample}
Prove by cases that $n(n+1)$ is even for every integer $n$.
\end{CoreProblem}
\begin{CoreProblem}{E.5}{Strategy choice and repair}
The following proof is intended to show that $A\subseteq B$ whenever
$A\cap B=A$. Identify the gap and repair it.

``Let $x\in A\cap B$. Then $x\in B$. Therefore $A\subseteq B$.''
\end{CoreProblem}
\end{multicols}
\CoreWorkArea
\FormBreak
\CoreHeader{Proofs}{P02}{A--E}{\FivePartScoring Use complete sentences, name the requested method or rule, and justify every inference.}
\begin{multicols}{2}
\begin{CoreProblem}{A.1}{Inference and proof structure}
Suppose that if the server is overloaded, requests are delayed; if requests are
delayed, the monitoring system sends an alert; and the server is overloaded.
Derive that the monitoring system sends an alert, naming the rule used at each
step.
\end{CoreProblem}
\begin{CoreProblem}{B.4}{Direct proof}
Give a direct proof that the difference of two even integers is even.
\end{CoreProblem}
\begin{CoreProblem}{C.3}{Contraposition or contradiction}
Prove by contradiction that there is no largest integer.
\end{CoreProblem}
\columnbreak
\begin{CoreProblem}{D.4}{Cases, existence, or counterexample}
Disprove each universal statement by giving a counterexample.
\begin{itemize}
  \item (a) The sum of two irrational numbers is irrational.
  \item (b) $n^2+n+41$ is prime for every nonnegative integer $n$.
\end{itemize}
\end{CoreProblem}
\begin{CoreProblem}{E.3}{Strategy choice and repair}
The following proof is intended to show that the product of a rational number
and an irrational number is irrational. Identify why the argument fails when
$r=0$. Add the weakest necessary restriction on $r$, state the corrected
theorem, and prove it.

``Let $r$ be rational and $x$ irrational. If $rx$ were rational, then $x=(rx)/r$ would be rational, a contradiction.''
\end{CoreProblem}
\end{multicols}
\CoreWorkArea
\FormBreak
\CoreHeader{Proofs}{P03}{A--E}{\FivePartScoring Use complete sentences, name the requested method or rule, and justify every inference.}
\begin{multicols}{2}
\begin{CoreProblem}{A.3}{Inference and proof structure}
Over one common domain, identify the error in the proposed inference. Then give
a two-element domain and interpretations of $P$ and $Q$ that make both premises
true and the conclusion false.
\[
\exists x\,P(x),\qquad \exists x\,Q(x)
\quad\therefore\quad
\exists x\,(P(x)\land Q(x)).
\]
\end{CoreProblem}
\begin{CoreProblem}{B.4}{Direct proof}
Give a direct proof that the difference of two even integers is even.
\end{CoreProblem}
\begin{CoreProblem}{C.1}{Contraposition or contradiction}
Let $m,n\in\mathbb Z$. Prove by contraposition that if $mn$ is odd, then both
$m$ and $n$ are odd. Your proof must state the contrapositive explicitly and
handle its disjunction by cases.
\end{CoreProblem}
\columnbreak
\begin{CoreProblem}{D.1}{Cases, existence, or counterexample}
Prove by cases that $n(n+1)$ is even for every integer $n$.
\end{CoreProblem}
\begin{CoreProblem}{E.1}{Strategy choice and repair}
For each statement, choose a promising first strategy---direct proof,
contraposition, contradiction, cases, constructive existence, or
counterexample---and state the first substantive step. Do not write the full
proofs.
\begin{itemize}
  \item (a) If $n^2$ is even, then the integer $n$ is even.
  \item (b) There is no smallest positive real number.
\end{itemize}
\end{CoreProblem}
\end{multicols}
\CoreWorkArea
\FormBreak
\CoreHeader{Proofs}{P04}{A--E}{\FivePartScoring Use complete sentences, name the requested method or rule, and justify every inference.}
\begin{multicols}{2}
\begin{CoreProblem}{A.1}{Inference and proof structure}
Suppose that if the server is overloaded, requests are delayed; if requests are
delayed, the monitoring system sends an alert; and the server is overloaded.
Derive that the monitoring system sends an alert, naming the rule used at each
step.
\end{CoreProblem}
\begin{CoreProblem}{B.3}{Direct proof}
Give a direct proof that the sum of two odd integers is even.
\end{CoreProblem}
\begin{CoreProblem}{C.4}{Contraposition or contradiction}
Let $r$ be rational and $x$ irrational. Prove by contradiction that $r+x$ is irrational.
\end{CoreProblem}
\columnbreak
\begin{CoreProblem}{D.4}{Cases, existence, or counterexample}
Disprove each universal statement by giving a counterexample.
\begin{itemize}
  \item (a) The sum of two irrational numbers is irrational.
  \item (b) $n^2+n+41$ is prime for every nonnegative integer $n$.
\end{itemize}
\end{CoreProblem}
\begin{CoreProblem}{E.3}{Strategy choice and repair}
The following proof is intended to show that the product of a rational number
and an irrational number is irrational. Identify why the argument fails when
$r=0$. Add the weakest necessary restriction on $r$, state the corrected
theorem, and prove it.

``Let $r$ be rational and $x$ irrational. If $rx$ were rational, then $x=(rx)/r$ would be rational, a contradiction.''
\end{CoreProblem}
\end{multicols}
\CoreWorkArea
\FormBreak
\CoreHeader{Proofs}{P05}{A--E}{\FivePartScoring Use complete sentences, name the requested method or rule, and justify every inference.}
\begin{multicols}{2}
\begin{CoreProblem}{A.3}{Inference and proof structure}
Over one common domain, identify the error in the proposed inference. Then give
a two-element domain and interpretations of $P$ and $Q$ that make both premises
true and the conclusion false.
\[
\exists x\,P(x),\qquad \exists x\,Q(x)
\quad\therefore\quad
\exists x\,(P(x)\land Q(x)).
\]
\end{CoreProblem}
\begin{CoreProblem}{B.5}{Direct proof}
Let $A,B,C$ be sets. Give a direct element proof that if $A\subseteq B$, then
$A\cap C\subseteq B\cap C$.
\end{CoreProblem}
\begin{CoreProblem}{C.5}{Contraposition or contradiction}
Prove by contraposition that if $n^2$ is even, then the integer $n$ is even.
\end{CoreProblem}
\columnbreak
\begin{CoreProblem}{D.1}{Cases, existence, or counterexample}
Prove by cases that $n(n+1)$ is even for every integer $n$.
\end{CoreProblem}
\begin{CoreProblem}{E.3}{Strategy choice and repair}
The following proof is intended to show that the product of a rational number
and an irrational number is irrational. Identify why the argument fails when
$r=0$. Add the weakest necessary restriction on $r$, state the corrected
theorem, and prove it.

``Let $r$ be rational and $x$ irrational. If $rx$ were rational, then $x=(rx)/r$ would be rational, a contradiction.''
\end{CoreProblem}
\end{multicols}
\CoreWorkArea
\FormBreak
\CoreHeader{Proofs}{P06}{A--E}{\FivePartScoring Use complete sentences, name the requested method or rule, and justify every inference.}
\begin{multicols}{2}
\begin{CoreProblem}{A.1}{Inference and proof structure}
Suppose that if the server is overloaded, requests are delayed; if requests are
delayed, the monitoring system sends an alert; and the server is overloaded.
Derive that the monitoring system sends an alert, naming the rule used at each
step.
\end{CoreProblem}
\begin{CoreProblem}{B.2}{Direct proof}
Give a direct proof that the sum of two rational numbers is rational.
\end{CoreProblem}
\begin{CoreProblem}{C.2}{Contraposition or contradiction}
Prove by contraposition that if $3n+2$ is odd, then the integer $n$ is odd.
\end{CoreProblem}
\columnbreak
\begin{CoreProblem}{D.1}{Cases, existence, or counterexample}
Prove by cases that $n(n+1)$ is even for every integer $n$.
\end{CoreProblem}
\begin{CoreProblem}{E.3}{Strategy choice and repair}
The following proof is intended to show that the product of a rational number
and an irrational number is irrational. Identify why the argument fails when
$r=0$. Add the weakest necessary restriction on $r$, state the corrected
theorem, and prove it.

``Let $r$ be rational and $x$ irrational. If $rx$ were rational, then $x=(rx)/r$ would be rational, a contradiction.''
\end{CoreProblem}
\end{multicols}
\CoreWorkArea
\FormBreak
\CoreHeader{Proofs}{P07}{A--E}{\FivePartScoring Use complete sentences, name the requested method or rule, and justify every inference.}
\begin{multicols}{2}
\begin{CoreProblem}{A.4}{Inference and proof structure}
From $p\to q$, $q\to r$, and $\neg r$, derive $\neg p$. Name the rule used at
each step.
\end{CoreProblem}
\begin{CoreProblem}{B.1}{Direct proof}
Give a direct proof that if an integer $n$ is even, then $n^2$ is even.
\end{CoreProblem}
\begin{CoreProblem}{C.1}{Contraposition or contradiction}
Let $m,n\in\mathbb Z$. Prove by contraposition that if $mn$ is odd, then both
$m$ and $n$ are odd. Your proof must state the contrapositive explicitly and
handle its disjunction by cases.
\end{CoreProblem}
\columnbreak
\begin{CoreProblem}{D.5}{Cases, existence, or counterexample}
Give a constructive existence proof that there are two irrational real numbers
whose product is rational.
\end{CoreProblem}
\begin{CoreProblem}{E.2}{Strategy choice and repair}
The following proof is intended to show that the sum of two even integers is even. Identify the gap and repair it.

``Let $a$ and $b$ be even. Then $a=2k$ and $b=2k$ for some integer $k$. Hence $a+b=4k$, which is even.''
\end{CoreProblem}
\end{multicols}
\CoreWorkArea
\FormBreak
\CoreHeader{Proofs}{P08}{A--E}{\FivePartScoring Use complete sentences, name the requested method or rule, and justify every inference.}
\begin{multicols}{2}
\begin{CoreProblem}{A.2}{Inference and proof structure}
Determine whether each argument is valid. Name the valid inference rule or the invalid form.
\begin{itemize}
  \item (a) If the program compiles, then it has no syntax errors. It has no syntax errors. Therefore, it compiles.
  \item (b) If the network is down, then the site is unavailable. The site is
        available. Therefore, the network is not down.
\end{itemize}
\end{CoreProblem}
\begin{CoreProblem}{B.3}{Direct proof}
Give a direct proof that the sum of two odd integers is even.
\end{CoreProblem}
\begin{CoreProblem}{C.3}{Contraposition or contradiction}
Prove by contradiction that there is no largest integer.
\end{CoreProblem}
\columnbreak
\begin{CoreProblem}{D.5}{Cases, existence, or counterexample}
Give a constructive existence proof that there are two irrational real numbers
whose product is rational.
\end{CoreProblem}
\begin{CoreProblem}{E.5}{Strategy choice and repair}
The following proof is intended to show that $A\subseteq B$ whenever
$A\cap B=A$. Identify the gap and repair it.

``Let $x\in A\cap B$. Then $x\in B$. Therefore $A\subseteq B$.''
\end{CoreProblem}
\end{multicols}
\CoreWorkArea
\FormBreak
\CoreHeader{Proofs}{P09}{A--E}{\FivePartScoring Use complete sentences, name the requested method or rule, and justify every inference.}
\begin{multicols}{2}
\begin{CoreProblem}{A.2}{Inference and proof structure}
Determine whether each argument is valid. Name the valid inference rule or the invalid form.
\begin{itemize}
  \item (a) If the program compiles, then it has no syntax errors. It has no syntax errors. Therefore, it compiles.
  \item (b) If the network is down, then the site is unavailable. The site is
        available. Therefore, the network is not down.
\end{itemize}
\end{CoreProblem}
\begin{CoreProblem}{B.1}{Direct proof}
Give a direct proof that if an integer $n$ is even, then $n^2$ is even.
\end{CoreProblem}
\begin{CoreProblem}{C.3}{Contraposition or contradiction}
Prove by contradiction that there is no largest integer.
\end{CoreProblem}
\columnbreak
\begin{CoreProblem}{D.2}{Cases, existence, or counterexample}
Prove by cases that $|x|\ge 0$ for every real number $x$, with equality exactly when $x=0$.
\end{CoreProblem}
\begin{CoreProblem}{E.2}{Strategy choice and repair}
The following proof is intended to show that the sum of two even integers is even. Identify the gap and repair it.

``Let $a$ and $b$ be even. Then $a=2k$ and $b=2k$ for some integer $k$. Hence $a+b=4k$, which is even.''
\end{CoreProblem}
\end{multicols}
\CoreWorkArea
\FormBreak
\CoreHeader{Proofs}{P10}{A--E}{\FivePartScoring Use complete sentences, name the requested method or rule, and justify every inference.}
\begin{multicols}{2}
\begin{CoreProblem}{A.4}{Inference and proof structure}
From $p\to q$, $q\to r$, and $\neg r$, derive $\neg p$. Name the rule used at
each step.
\end{CoreProblem}
\begin{CoreProblem}{B.1}{Direct proof}
Give a direct proof that if an integer $n$ is even, then $n^2$ is even.
\end{CoreProblem}
\begin{CoreProblem}{C.2}{Contraposition or contradiction}
Prove by contraposition that if $3n+2$ is odd, then the integer $n$ is odd.
\end{CoreProblem}
\columnbreak
\begin{CoreProblem}{D.2}{Cases, existence, or counterexample}
Prove by cases that $|x|\ge 0$ for every real number $x$, with equality exactly when $x=0$.
\end{CoreProblem}
\begin{CoreProblem}{E.5}{Strategy choice and repair}
The following proof is intended to show that $A\subseteq B$ whenever
$A\cap B=A$. Identify the gap and repair it.

``Let $x\in A\cap B$. Then $x\in B$. Therefore $A\subseteq B$.''
\end{CoreProblem}
\end{multicols}
\CoreWorkArea
\FormBreak
\CoreHeader{Proofs}{P11}{A--E}{\FivePartScoring Use complete sentences, name the requested method or rule, and justify every inference.}
\begin{multicols}{2}
\begin{CoreProblem}{A.5}{Inference and proof structure}
Over a nonempty domain, use the premises
$\forall x\,(P(x)\to Q(x))$ and $\forall x\,P(x)$ to derive
$\forall x\,Q(x)$. Name the rule used at each step.
\end{CoreProblem}
\begin{CoreProblem}{B.3}{Direct proof}
Give a direct proof that the sum of two odd integers is even.
\end{CoreProblem}
\begin{CoreProblem}{C.2}{Contraposition or contradiction}
Prove by contraposition that if $3n+2$ is odd, then the integer $n$ is odd.
\end{CoreProblem}
\columnbreak
\begin{CoreProblem}{D.2}{Cases, existence, or counterexample}
Prove by cases that $|x|\ge 0$ for every real number $x$, with equality exactly when $x=0$.
\end{CoreProblem}
\begin{CoreProblem}{E.4}{Strategy choice and repair}
The following argument is circular. Explain the circularity, then replace it with a valid proof.

``To prove that $n^2$ is odd whenever $n$ is odd, assume $n^2$ is odd. Thus
$n^2=2k+1$ for some integer $k$, so $n^2$ is odd.''
\end{CoreProblem}
\end{multicols}
\CoreWorkArea
\FormBreak
\CoreHeader{Proofs}{P12}{A--E}{\FivePartScoring Use complete sentences, name the requested method or rule, and justify every inference.}
\begin{multicols}{2}
\begin{CoreProblem}{A.4}{Inference and proof structure}
From $p\to q$, $q\to r$, and $\neg r$, derive $\neg p$. Name the rule used at
each step.
\end{CoreProblem}
\begin{CoreProblem}{B.1}{Direct proof}
Give a direct proof that if an integer $n$ is even, then $n^2$ is even.
\end{CoreProblem}
\begin{CoreProblem}{C.1}{Contraposition or contradiction}
Let $m,n\in\mathbb Z$. Prove by contraposition that if $mn$ is odd, then both
$m$ and $n$ are odd. Your proof must state the contrapositive explicitly and
handle its disjunction by cases.
\end{CoreProblem}
\columnbreak
\begin{CoreProblem}{D.4}{Cases, existence, or counterexample}
Disprove each universal statement by giving a counterexample.
\begin{itemize}
  \item (a) The sum of two irrational numbers is irrational.
  \item (b) $n^2+n+41$ is prime for every nonnegative integer $n$.
\end{itemize}
\end{CoreProblem}
\begin{CoreProblem}{E.1}{Strategy choice and repair}
For each statement, choose a promising first strategy---direct proof,
contraposition, contradiction, cases, constructive existence, or
counterexample---and state the first substantive step. Do not write the full
proofs.
\begin{itemize}
  \item (a) If $n^2$ is even, then the integer $n$ is even.
  \item (b) There is no smallest positive real number.
\end{itemize}
\end{CoreProblem}
\end{multicols}
\CoreWorkArea
\FormBreak
\CoreHeader{Proofs}{P13}{A--E}{\FivePartScoring Use complete sentences, name the requested method or rule, and justify every inference.}
\begin{multicols}{2}
\begin{CoreProblem}{A.3}{Inference and proof structure}
Over one common domain, identify the error in the proposed inference. Then give
a two-element domain and interpretations of $P$ and $Q$ that make both premises
true and the conclusion false.
\[
\exists x\,P(x),\qquad \exists x\,Q(x)
\quad\therefore\quad
\exists x\,(P(x)\land Q(x)).
\]
\end{CoreProblem}
\begin{CoreProblem}{B.3}{Direct proof}
Give a direct proof that the sum of two odd integers is even.
\end{CoreProblem}
\begin{CoreProblem}{C.4}{Contraposition or contradiction}
Let $r$ be rational and $x$ irrational. Prove by contradiction that $r+x$ is irrational.
\end{CoreProblem}
\columnbreak
\begin{CoreProblem}{D.4}{Cases, existence, or counterexample}
Disprove each universal statement by giving a counterexample.
\begin{itemize}
  \item (a) The sum of two irrational numbers is irrational.
  \item (b) $n^2+n+41$ is prime for every nonnegative integer $n$.
\end{itemize}
\end{CoreProblem}
\begin{CoreProblem}{E.4}{Strategy choice and repair}
The following argument is circular. Explain the circularity, then replace it with a valid proof.

``To prove that $n^2$ is odd whenever $n$ is odd, assume $n^2$ is odd. Thus
$n^2=2k+1$ for some integer $k$, so $n^2$ is odd.''
\end{CoreProblem}
\end{multicols}
\CoreWorkArea
\FormBreak
\CoreHeader{Proofs}{P14}{A--E}{\FivePartScoring Use complete sentences, name the requested method or rule, and justify every inference.}
\begin{multicols}{2}
\begin{CoreProblem}{A.2}{Inference and proof structure}
Determine whether each argument is valid. Name the valid inference rule or the invalid form.
\begin{itemize}
  \item (a) If the program compiles, then it has no syntax errors. It has no syntax errors. Therefore, it compiles.
  \item (b) If the network is down, then the site is unavailable. The site is
        available. Therefore, the network is not down.
\end{itemize}
\end{CoreProblem}
\begin{CoreProblem}{B.3}{Direct proof}
Give a direct proof that the sum of two odd integers is even.
\end{CoreProblem}
\begin{CoreProblem}{C.3}{Contraposition or contradiction}
Prove by contradiction that there is no largest integer.
\end{CoreProblem}
\columnbreak
\begin{CoreProblem}{D.2}{Cases, existence, or counterexample}
Prove by cases that $|x|\ge 0$ for every real number $x$, with equality exactly when $x=0$.
\end{CoreProblem}
\begin{CoreProblem}{E.4}{Strategy choice and repair}
The following argument is circular. Explain the circularity, then replace it with a valid proof.

``To prove that $n^2$ is odd whenever $n$ is odd, assume $n^2$ is odd. Thus
$n^2=2k+1$ for some integer $k$, so $n^2$ is odd.''
\end{CoreProblem}
\end{multicols}
\CoreWorkArea
\FormBreak
\CoreHeader{Proofs}{P15}{A--E}{\FivePartScoring Use complete sentences, name the requested method or rule, and justify every inference.}
\begin{multicols}{2}
\begin{CoreProblem}{A.4}{Inference and proof structure}
From $p\to q$, $q\to r$, and $\neg r$, derive $\neg p$. Name the rule used at
each step.
\end{CoreProblem}
\begin{CoreProblem}{B.5}{Direct proof}
Let $A,B,C$ be sets. Give a direct element proof that if $A\subseteq B$, then
$A\cap C\subseteq B\cap C$.
\end{CoreProblem}
\begin{CoreProblem}{C.3}{Contraposition or contradiction}
Prove by contradiction that there is no largest integer.
\end{CoreProblem}
\columnbreak
\begin{CoreProblem}{D.4}{Cases, existence, or counterexample}
Disprove each universal statement by giving a counterexample.
\begin{itemize}
  \item (a) The sum of two irrational numbers is irrational.
  \item (b) $n^2+n+41$ is prime for every nonnegative integer $n$.
\end{itemize}
\end{CoreProblem}
\begin{CoreProblem}{E.5}{Strategy choice and repair}
The following proof is intended to show that $A\subseteq B$ whenever
$A\cap B=A$. Identify the gap and repair it.

``Let $x\in A\cap B$. Then $x\in B$. Therefore $A\subseteq B$.''
\end{CoreProblem}
\end{multicols}
\CoreWorkArea
\FormBreak
\CoreHeader{Proofs}{P16}{A--E}{\FivePartScoring Use complete sentences, name the requested method or rule, and justify every inference.}
\begin{multicols}{2}
\begin{CoreProblem}{A.5}{Inference and proof structure}
Over a nonempty domain, use the premises
$\forall x\,(P(x)\to Q(x))$ and $\forall x\,P(x)$ to derive
$\forall x\,Q(x)$. Name the rule used at each step.
\end{CoreProblem}
\begin{CoreProblem}{B.4}{Direct proof}
Give a direct proof that the difference of two even integers is even.
\end{CoreProblem}
\begin{CoreProblem}{C.4}{Contraposition or contradiction}
Let $r$ be rational and $x$ irrational. Prove by contradiction that $r+x$ is irrational.
\end{CoreProblem}
\columnbreak
\begin{CoreProblem}{D.4}{Cases, existence, or counterexample}
Disprove each universal statement by giving a counterexample.
\begin{itemize}
  \item (a) The sum of two irrational numbers is irrational.
  \item (b) $n^2+n+41$ is prime for every nonnegative integer $n$.
\end{itemize}
\end{CoreProblem}
\begin{CoreProblem}{E.5}{Strategy choice and repair}
The following proof is intended to show that $A\subseteq B$ whenever
$A\cap B=A$. Identify the gap and repair it.

``Let $x\in A\cap B$. Then $x\in B$. Therefore $A\subseteq B$.''
\end{CoreProblem}
\end{multicols}
\CoreWorkArea
\FormBreak
\CoreHeader{Proofs}{P17}{A--E}{\FivePartScoring Use complete sentences, name the requested method or rule, and justify every inference.}
\begin{multicols}{2}
\begin{CoreProblem}{A.3}{Inference and proof structure}
Over one common domain, identify the error in the proposed inference. Then give
a two-element domain and interpretations of $P$ and $Q$ that make both premises
true and the conclusion false.
\[
\exists x\,P(x),\qquad \exists x\,Q(x)
\quad\therefore\quad
\exists x\,(P(x)\land Q(x)).
\]
\end{CoreProblem}
\begin{CoreProblem}{B.5}{Direct proof}
Let $A,B,C$ be sets. Give a direct element proof that if $A\subseteq B$, then
$A\cap C\subseteq B\cap C$.
\end{CoreProblem}
\begin{CoreProblem}{C.5}{Contraposition or contradiction}
Prove by contraposition that if $n^2$ is even, then the integer $n$ is even.
\end{CoreProblem}
\columnbreak
\begin{CoreProblem}{D.4}{Cases, existence, or counterexample}
Disprove each universal statement by giving a counterexample.
\begin{itemize}
  \item (a) The sum of two irrational numbers is irrational.
  \item (b) $n^2+n+41$ is prime for every nonnegative integer $n$.
\end{itemize}
\end{CoreProblem}
\begin{CoreProblem}{E.2}{Strategy choice and repair}
The following proof is intended to show that the sum of two even integers is even. Identify the gap and repair it.

``Let $a$ and $b$ be even. Then $a=2k$ and $b=2k$ for some integer $k$. Hence $a+b=4k$, which is even.''
\end{CoreProblem}
\end{multicols}
\CoreWorkArea
\FormBreak
\CoreHeader{Proofs}{P18}{A--E}{\FivePartScoring Use complete sentences, name the requested method or rule, and justify every inference.}
\begin{multicols}{2}
\begin{CoreProblem}{A.1}{Inference and proof structure}
Suppose that if the server is overloaded, requests are delayed; if requests are
delayed, the monitoring system sends an alert; and the server is overloaded.
Derive that the monitoring system sends an alert, naming the rule used at each
step.
\end{CoreProblem}
\begin{CoreProblem}{B.4}{Direct proof}
Give a direct proof that the difference of two even integers is even.
\end{CoreProblem}
\begin{CoreProblem}{C.2}{Contraposition or contradiction}
Prove by contraposition that if $3n+2$ is odd, then the integer $n$ is odd.
\end{CoreProblem}
\columnbreak
\begin{CoreProblem}{D.3}{Cases, existence, or counterexample}
Give a constructive existence proof that there are four consecutive composite
positive integers. Verify explicitly that every integer in your example is
composite.
\end{CoreProblem}
\begin{CoreProblem}{E.4}{Strategy choice and repair}
The following argument is circular. Explain the circularity, then replace it with a valid proof.

``To prove that $n^2$ is odd whenever $n$ is odd, assume $n^2$ is odd. Thus
$n^2=2k+1$ for some integer $k$, so $n^2$ is odd.''
\end{CoreProblem}
\end{multicols}
\CoreWorkArea
\FormBreak
\CoreHeader{Proofs}{P19}{A--E}{\FivePartScoring Use complete sentences, name the requested method or rule, and justify every inference.}
\begin{multicols}{2}
\begin{CoreProblem}{A.4}{Inference and proof structure}
From $p\to q$, $q\to r$, and $\neg r$, derive $\neg p$. Name the rule used at
each step.
\end{CoreProblem}
\begin{CoreProblem}{B.4}{Direct proof}
Give a direct proof that the difference of two even integers is even.
\end{CoreProblem}
\begin{CoreProblem}{C.5}{Contraposition or contradiction}
Prove by contraposition that if $n^2$ is even, then the integer $n$ is even.
\end{CoreProblem}
\columnbreak
\begin{CoreProblem}{D.2}{Cases, existence, or counterexample}
Prove by cases that $|x|\ge 0$ for every real number $x$, with equality exactly when $x=0$.
\end{CoreProblem}
\begin{CoreProblem}{E.5}{Strategy choice and repair}
The following proof is intended to show that $A\subseteq B$ whenever
$A\cap B=A$. Identify the gap and repair it.

``Let $x\in A\cap B$. Then $x\in B$. Therefore $A\subseteq B$.''
\end{CoreProblem}
\end{multicols}
\CoreWorkArea
\FormBreak
\CoreHeader{Proofs}{P20}{A--E}{\FivePartScoring Use complete sentences, name the requested method or rule, and justify every inference.}
\begin{multicols}{2}
\begin{CoreProblem}{A.3}{Inference and proof structure}
Over one common domain, identify the error in the proposed inference. Then give
a two-element domain and interpretations of $P$ and $Q$ that make both premises
true and the conclusion false.
\[
\exists x\,P(x),\qquad \exists x\,Q(x)
\quad\therefore\quad
\exists x\,(P(x)\land Q(x)).
\]
\end{CoreProblem}
\begin{CoreProblem}{B.1}{Direct proof}
Give a direct proof that if an integer $n$ is even, then $n^2$ is even.
\end{CoreProblem}
\begin{CoreProblem}{C.5}{Contraposition or contradiction}
Prove by contraposition that if $n^2$ is even, then the integer $n$ is even.
\end{CoreProblem}
\columnbreak
\begin{CoreProblem}{D.2}{Cases, existence, or counterexample}
Prove by cases that $|x|\ge 0$ for every real number $x$, with equality exactly when $x=0$.
\end{CoreProblem}
\begin{CoreProblem}{E.4}{Strategy choice and repair}
The following argument is circular. Explain the circularity, then replace it with a valid proof.

``To prove that $n^2$ is odd whenever $n$ is odd, assume $n^2$ is odd. Thus
$n^2=2k+1$ for some integer $k$, so $n^2$ is odd.''
\end{CoreProblem}
\end{multicols}
\CoreWorkArea
\FormBreak
\CoreHeader{Proofs}{P21}{A--E}{\FivePartScoring Use complete sentences, name the requested method or rule, and justify every inference.}
\begin{multicols}{2}
\begin{CoreProblem}{A.1}{Inference and proof structure}
Suppose that if the server is overloaded, requests are delayed; if requests are
delayed, the monitoring system sends an alert; and the server is overloaded.
Derive that the monitoring system sends an alert, naming the rule used at each
step.
\end{CoreProblem}
\begin{CoreProblem}{B.4}{Direct proof}
Give a direct proof that the difference of two even integers is even.
\end{CoreProblem}
\begin{CoreProblem}{C.2}{Contraposition or contradiction}
Prove by contraposition that if $3n+2$ is odd, then the integer $n$ is odd.
\end{CoreProblem}
\columnbreak
\begin{CoreProblem}{D.1}{Cases, existence, or counterexample}
Prove by cases that $n(n+1)$ is even for every integer $n$.
\end{CoreProblem}
\begin{CoreProblem}{E.5}{Strategy choice and repair}
The following proof is intended to show that $A\subseteq B$ whenever
$A\cap B=A$. Identify the gap and repair it.

``Let $x\in A\cap B$. Then $x\in B$. Therefore $A\subseteq B$.''
\end{CoreProblem}
\end{multicols}
\CoreWorkArea
\FormBreak
\CoreHeader{Proofs}{P22}{A--E}{\FivePartScoring Use complete sentences, name the requested method or rule, and justify every inference.}
\begin{multicols}{2}
\begin{CoreProblem}{A.2}{Inference and proof structure}
Determine whether each argument is valid. Name the valid inference rule or the invalid form.
\begin{itemize}
  \item (a) If the program compiles, then it has no syntax errors. It has no syntax errors. Therefore, it compiles.
  \item (b) If the network is down, then the site is unavailable. The site is
        available. Therefore, the network is not down.
\end{itemize}
\end{CoreProblem}
\begin{CoreProblem}{B.4}{Direct proof}
Give a direct proof that the difference of two even integers is even.
\end{CoreProblem}
\begin{CoreProblem}{C.5}{Contraposition or contradiction}
Prove by contraposition that if $n^2$ is even, then the integer $n$ is even.
\end{CoreProblem}
\columnbreak
\begin{CoreProblem}{D.2}{Cases, existence, or counterexample}
Prove by cases that $|x|\ge 0$ for every real number $x$, with equality exactly when $x=0$.
\end{CoreProblem}
\begin{CoreProblem}{E.1}{Strategy choice and repair}
For each statement, choose a promising first strategy---direct proof,
contraposition, contradiction, cases, constructive existence, or
counterexample---and state the first substantive step. Do not write the full
proofs.
\begin{itemize}
  \item (a) If $n^2$ is even, then the integer $n$ is even.
  \item (b) There is no smallest positive real number.
\end{itemize}
\end{CoreProblem}
\end{multicols}
\CoreWorkArea
\FormBreak
\CoreHeader{Proofs}{P23}{A--E}{\FivePartScoring Use complete sentences, name the requested method or rule, and justify every inference.}
\begin{multicols}{2}
\begin{CoreProblem}{A.2}{Inference and proof structure}
Determine whether each argument is valid. Name the valid inference rule or the invalid form.
\begin{itemize}
  \item (a) If the program compiles, then it has no syntax errors. It has no syntax errors. Therefore, it compiles.
  \item (b) If the network is down, then the site is unavailable. The site is
        available. Therefore, the network is not down.
\end{itemize}
\end{CoreProblem}
\begin{CoreProblem}{B.4}{Direct proof}
Give a direct proof that the difference of two even integers is even.
\end{CoreProblem}
\begin{CoreProblem}{C.3}{Contraposition or contradiction}
Prove by contradiction that there is no largest integer.
\end{CoreProblem}
\columnbreak
\begin{CoreProblem}{D.3}{Cases, existence, or counterexample}
Give a constructive existence proof that there are four consecutive composite
positive integers. Verify explicitly that every integer in your example is
composite.
\end{CoreProblem}
\begin{CoreProblem}{E.3}{Strategy choice and repair}
The following proof is intended to show that the product of a rational number
and an irrational number is irrational. Identify why the argument fails when
$r=0$. Add the weakest necessary restriction on $r$, state the corrected
theorem, and prove it.

``Let $r$ be rational and $x$ irrational. If $rx$ were rational, then $x=(rx)/r$ would be rational, a contradiction.''
\end{CoreProblem}
\end{multicols}
\CoreWorkArea
\FormBreak
\CoreHeader{Proofs}{P24}{A--E}{\FivePartScoring Use complete sentences, name the requested method or rule, and justify every inference.}
\begin{multicols}{2}
\begin{CoreProblem}{A.2}{Inference and proof structure}
Determine whether each argument is valid. Name the valid inference rule or the invalid form.
\begin{itemize}
  \item (a) If the program compiles, then it has no syntax errors. It has no syntax errors. Therefore, it compiles.
  \item (b) If the network is down, then the site is unavailable. The site is
        available. Therefore, the network is not down.
\end{itemize}
\end{CoreProblem}
\begin{CoreProblem}{B.2}{Direct proof}
Give a direct proof that the sum of two rational numbers is rational.
\end{CoreProblem}
\begin{CoreProblem}{C.3}{Contraposition or contradiction}
Prove by contradiction that there is no largest integer.
\end{CoreProblem}
\columnbreak
\begin{CoreProblem}{D.1}{Cases, existence, or counterexample}
Prove by cases that $n(n+1)$ is even for every integer $n$.
\end{CoreProblem}
\begin{CoreProblem}{E.1}{Strategy choice and repair}
For each statement, choose a promising first strategy---direct proof,
contraposition, contradiction, cases, constructive existence, or
counterexample---and state the first substantive step. Do not write the full
proofs.
\begin{itemize}
  \item (a) If $n^2$ is even, then the integer $n$ is even.
  \item (b) There is no smallest positive real number.
\end{itemize}
\end{CoreProblem}
\end{multicols}
\CoreWorkArea
\FormBreak
\CoreHeader{Proofs}{P25}{A--E}{\FivePartScoring Use complete sentences, name the requested method or rule, and justify every inference.}
\begin{multicols}{2}
\begin{CoreProblem}{A.4}{Inference and proof structure}
From $p\to q$, $q\to r$, and $\neg r$, derive $\neg p$. Name the rule used at
each step.
\end{CoreProblem}
\begin{CoreProblem}{B.2}{Direct proof}
Give a direct proof that the sum of two rational numbers is rational.
\end{CoreProblem}
\begin{CoreProblem}{C.2}{Contraposition or contradiction}
Prove by contraposition that if $3n+2$ is odd, then the integer $n$ is odd.
\end{CoreProblem}
\columnbreak
\begin{CoreProblem}{D.5}{Cases, existence, or counterexample}
Give a constructive existence proof that there are two irrational real numbers
whose product is rational.
\end{CoreProblem}
\begin{CoreProblem}{E.3}{Strategy choice and repair}
The following proof is intended to show that the product of a rational number
and an irrational number is irrational. Identify why the argument fails when
$r=0$. Add the weakest necessary restriction on $r$, state the corrected
theorem, and prove it.

``Let $r$ be rational and $x$ irrational. If $rx$ were rational, then $x=(rx)/r$ would be rational, a contradiction.''
\end{CoreProblem}
\end{multicols}
\CoreWorkArea
\FormBreak
\CoreHeader{Proofs}{P26}{A--E}{\FivePartScoring Use complete sentences, name the requested method or rule, and justify every inference.}
\begin{multicols}{2}
\begin{CoreProblem}{A.5}{Inference and proof structure}
Over a nonempty domain, use the premises
$\forall x\,(P(x)\to Q(x))$ and $\forall x\,P(x)$ to derive
$\forall x\,Q(x)$. Name the rule used at each step.
\end{CoreProblem}
\begin{CoreProblem}{B.3}{Direct proof}
Give a direct proof that the sum of two odd integers is even.
\end{CoreProblem}
\begin{CoreProblem}{C.4}{Contraposition or contradiction}
Let $r$ be rational and $x$ irrational. Prove by contradiction that $r+x$ is irrational.
\end{CoreProblem}
\columnbreak
\begin{CoreProblem}{D.3}{Cases, existence, or counterexample}
Give a constructive existence proof that there are four consecutive composite
positive integers. Verify explicitly that every integer in your example is
composite.
\end{CoreProblem}
\begin{CoreProblem}{E.4}{Strategy choice and repair}
The following argument is circular. Explain the circularity, then replace it with a valid proof.

``To prove that $n^2$ is odd whenever $n$ is odd, assume $n^2$ is odd. Thus
$n^2=2k+1$ for some integer $k$, so $n^2$ is odd.''
\end{CoreProblem}
\end{multicols}
\CoreWorkArea
\FormBreak
\CoreHeader{Proofs}{P27}{A--E}{\FivePartScoring Use complete sentences, name the requested method or rule, and justify every inference.}
\begin{multicols}{2}
\begin{CoreProblem}{A.5}{Inference and proof structure}
Over a nonempty domain, use the premises
$\forall x\,(P(x)\to Q(x))$ and $\forall x\,P(x)$ to derive
$\forall x\,Q(x)$. Name the rule used at each step.
\end{CoreProblem}
\begin{CoreProblem}{B.1}{Direct proof}
Give a direct proof that if an integer $n$ is even, then $n^2$ is even.
\end{CoreProblem}
\begin{CoreProblem}{C.3}{Contraposition or contradiction}
Prove by contradiction that there is no largest integer.
\end{CoreProblem}
\columnbreak
\begin{CoreProblem}{D.1}{Cases, existence, or counterexample}
Prove by cases that $n(n+1)$ is even for every integer $n$.
\end{CoreProblem}
\begin{CoreProblem}{E.4}{Strategy choice and repair}
The following argument is circular. Explain the circularity, then replace it with a valid proof.

``To prove that $n^2$ is odd whenever $n$ is odd, assume $n^2$ is odd. Thus
$n^2=2k+1$ for some integer $k$, so $n^2$ is odd.''
\end{CoreProblem}
\end{multicols}
\CoreWorkArea
\FormBreak
\CoreHeader{Proofs}{P28}{A--E}{\FivePartScoring Use complete sentences, name the requested method or rule, and justify every inference.}
\begin{multicols}{2}
\begin{CoreProblem}{A.4}{Inference and proof structure}
From $p\to q$, $q\to r$, and $\neg r$, derive $\neg p$. Name the rule used at
each step.
\end{CoreProblem}
\begin{CoreProblem}{B.5}{Direct proof}
Let $A,B,C$ be sets. Give a direct element proof that if $A\subseteq B$, then
$A\cap C\subseteq B\cap C$.
\end{CoreProblem}
\begin{CoreProblem}{C.1}{Contraposition or contradiction}
Let $m,n\in\mathbb Z$. Prove by contraposition that if $mn$ is odd, then both
$m$ and $n$ are odd. Your proof must state the contrapositive explicitly and
handle its disjunction by cases.
\end{CoreProblem}
\columnbreak
\begin{CoreProblem}{D.5}{Cases, existence, or counterexample}
Give a constructive existence proof that there are two irrational real numbers
whose product is rational.
\end{CoreProblem}
\begin{CoreProblem}{E.5}{Strategy choice and repair}
The following proof is intended to show that $A\subseteq B$ whenever
$A\cap B=A$. Identify the gap and repair it.

``Let $x\in A\cap B$. Then $x\in B$. Therefore $A\subseteq B$.''
\end{CoreProblem}
\end{multicols}
\CoreWorkArea
\FormBreak
\CoreHeader{Proofs}{P29}{A--E}{\FivePartScoring Use complete sentences, name the requested method or rule, and justify every inference.}
\begin{multicols}{2}
\begin{CoreProblem}{A.3}{Inference and proof structure}
Over one common domain, identify the error in the proposed inference. Then give
a two-element domain and interpretations of $P$ and $Q$ that make both premises
true and the conclusion false.
\[
\exists x\,P(x),\qquad \exists x\,Q(x)
\quad\therefore\quad
\exists x\,(P(x)\land Q(x)).
\]
\end{CoreProblem}
\begin{CoreProblem}{B.4}{Direct proof}
Give a direct proof that the difference of two even integers is even.
\end{CoreProblem}
\begin{CoreProblem}{C.5}{Contraposition or contradiction}
Prove by contraposition that if $n^2$ is even, then the integer $n$ is even.
\end{CoreProblem}
\columnbreak
\begin{CoreProblem}{D.4}{Cases, existence, or counterexample}
Disprove each universal statement by giving a counterexample.
\begin{itemize}
  \item (a) The sum of two irrational numbers is irrational.
  \item (b) $n^2+n+41$ is prime for every nonnegative integer $n$.
\end{itemize}
\end{CoreProblem}
\begin{CoreProblem}{E.1}{Strategy choice and repair}
For each statement, choose a promising first strategy---direct proof,
contraposition, contradiction, cases, constructive existence, or
counterexample---and state the first substantive step. Do not write the full
proofs.
\begin{itemize}
  \item (a) If $n^2$ is even, then the integer $n$ is even.
  \item (b) There is no smallest positive real number.
\end{itemize}
\end{CoreProblem}
\end{multicols}
\CoreWorkArea
\FormBreak
\CoreHeader{Proofs}{P30}{A--E}{\FivePartScoring Use complete sentences, name the requested method or rule, and justify every inference.}
\begin{multicols}{2}
\begin{CoreProblem}{A.2}{Inference and proof structure}
Determine whether each argument is valid. Name the valid inference rule or the invalid form.
\begin{itemize}
  \item (a) If the program compiles, then it has no syntax errors. It has no syntax errors. Therefore, it compiles.
  \item (b) If the network is down, then the site is unavailable. The site is
        available. Therefore, the network is not down.
\end{itemize}
\end{CoreProblem}
\begin{CoreProblem}{B.3}{Direct proof}
Give a direct proof that the sum of two odd integers is even.
\end{CoreProblem}
\begin{CoreProblem}{C.3}{Contraposition or contradiction}
Prove by contradiction that there is no largest integer.
\end{CoreProblem}
\columnbreak
\begin{CoreProblem}{D.2}{Cases, existence, or counterexample}
Prove by cases that $|x|\ge 0$ for every real number $x$, with equality exactly when $x=0$.
\end{CoreProblem}
\begin{CoreProblem}{E.3}{Strategy choice and repair}
The following proof is intended to show that the product of a rational number
and an irrational number is irrational. Identify why the argument fails when
$r=0$. Add the weakest necessary restriction on $r$, state the corrected
theorem, and prove it.

``Let $r$ be rational and $x$ irrational. If $rx$ were rational, then $x=(rx)/r$ would be rational, a contradiction.''
\end{CoreProblem}
\end{multicols}
\CoreWorkArea
\FormBreak
\CoreHeader{Proofs}{P31}{A--E}{\FivePartScoring Use complete sentences, name the requested method or rule, and justify every inference.}
\begin{multicols}{2}
\begin{CoreProblem}{A.1}{Inference and proof structure}
Suppose that if the server is overloaded, requests are delayed; if requests are
delayed, the monitoring system sends an alert; and the server is overloaded.
Derive that the monitoring system sends an alert, naming the rule used at each
step.
\end{CoreProblem}
\begin{CoreProblem}{B.3}{Direct proof}
Give a direct proof that the sum of two odd integers is even.
\end{CoreProblem}
\begin{CoreProblem}{C.5}{Contraposition or contradiction}
Prove by contraposition that if $n^2$ is even, then the integer $n$ is even.
\end{CoreProblem}
\columnbreak
\begin{CoreProblem}{D.5}{Cases, existence, or counterexample}
Give a constructive existence proof that there are two irrational real numbers
whose product is rational.
\end{CoreProblem}
\begin{CoreProblem}{E.2}{Strategy choice and repair}
The following proof is intended to show that the sum of two even integers is even. Identify the gap and repair it.

``Let $a$ and $b$ be even. Then $a=2k$ and $b=2k$ for some integer $k$. Hence $a+b=4k$, which is even.''
\end{CoreProblem}
\end{multicols}
\CoreWorkArea
\FormBreak
\CoreHeader{Proofs}{P32}{A--E}{\FivePartScoring Use complete sentences, name the requested method or rule, and justify every inference.}
\begin{multicols}{2}
\begin{CoreProblem}{A.1}{Inference and proof structure}
Suppose that if the server is overloaded, requests are delayed; if requests are
delayed, the monitoring system sends an alert; and the server is overloaded.
Derive that the monitoring system sends an alert, naming the rule used at each
step.
\end{CoreProblem}
\begin{CoreProblem}{B.1}{Direct proof}
Give a direct proof that if an integer $n$ is even, then $n^2$ is even.
\end{CoreProblem}
\begin{CoreProblem}{C.3}{Contraposition or contradiction}
Prove by contradiction that there is no largest integer.
\end{CoreProblem}
\columnbreak
\begin{CoreProblem}{D.1}{Cases, existence, or counterexample}
Prove by cases that $n(n+1)$ is even for every integer $n$.
\end{CoreProblem}
\begin{CoreProblem}{E.2}{Strategy choice and repair}
The following proof is intended to show that the sum of two even integers is even. Identify the gap and repair it.

``Let $a$ and $b$ be even. Then $a=2k$ and $b=2k$ for some integer $k$. Hence $a+b=4k$, which is even.''
\end{CoreProblem}
\end{multicols}
\CoreWorkArea
\FormBreak
\CoreHeader{Proofs}{P33}{A--E}{\FivePartScoring Use complete sentences, name the requested method or rule, and justify every inference.}
\begin{multicols}{2}
\begin{CoreProblem}{A.4}{Inference and proof structure}
From $p\to q$, $q\to r$, and $\neg r$, derive $\neg p$. Name the rule used at
each step.
\end{CoreProblem}
\begin{CoreProblem}{B.1}{Direct proof}
Give a direct proof that if an integer $n$ is even, then $n^2$ is even.
\end{CoreProblem}
\begin{CoreProblem}{C.1}{Contraposition or contradiction}
Let $m,n\in\mathbb Z$. Prove by contraposition that if $mn$ is odd, then both
$m$ and $n$ are odd. Your proof must state the contrapositive explicitly and
handle its disjunction by cases.
\end{CoreProblem}
\columnbreak
\begin{CoreProblem}{D.5}{Cases, existence, or counterexample}
Give a constructive existence proof that there are two irrational real numbers
whose product is rational.
\end{CoreProblem}
\begin{CoreProblem}{E.4}{Strategy choice and repair}
The following argument is circular. Explain the circularity, then replace it with a valid proof.

``To prove that $n^2$ is odd whenever $n$ is odd, assume $n^2$ is odd. Thus
$n^2=2k+1$ for some integer $k$, so $n^2$ is odd.''
\end{CoreProblem}
\end{multicols}
\CoreWorkArea
\FormBreak
\CoreHeader{Proofs}{P34}{A--E}{\FivePartScoring Use complete sentences, name the requested method or rule, and justify every inference.}
\begin{multicols}{2}
\begin{CoreProblem}{A.2}{Inference and proof structure}
Determine whether each argument is valid. Name the valid inference rule or the invalid form.
\begin{itemize}
  \item (a) If the program compiles, then it has no syntax errors. It has no syntax errors. Therefore, it compiles.
  \item (b) If the network is down, then the site is unavailable. The site is
        available. Therefore, the network is not down.
\end{itemize}
\end{CoreProblem}
\begin{CoreProblem}{B.3}{Direct proof}
Give a direct proof that the sum of two odd integers is even.
\end{CoreProblem}
\begin{CoreProblem}{C.4}{Contraposition or contradiction}
Let $r$ be rational and $x$ irrational. Prove by contradiction that $r+x$ is irrational.
\end{CoreProblem}
\columnbreak
\begin{CoreProblem}{D.4}{Cases, existence, or counterexample}
Disprove each universal statement by giving a counterexample.
\begin{itemize}
  \item (a) The sum of two irrational numbers is irrational.
  \item (b) $n^2+n+41$ is prime for every nonnegative integer $n$.
\end{itemize}
\end{CoreProblem}
\begin{CoreProblem}{E.2}{Strategy choice and repair}
The following proof is intended to show that the sum of two even integers is even. Identify the gap and repair it.

``Let $a$ and $b$ be even. Then $a=2k$ and $b=2k$ for some integer $k$. Hence $a+b=4k$, which is even.''
\end{CoreProblem}
\end{multicols}
\CoreWorkArea
\FormBreak
\CoreHeader{Proofs}{P35}{A--E}{\FivePartScoring Use complete sentences, name the requested method or rule, and justify every inference.}
\begin{multicols}{2}
\begin{CoreProblem}{A.2}{Inference and proof structure}
Determine whether each argument is valid. Name the valid inference rule or the invalid form.
\begin{itemize}
  \item (a) If the program compiles, then it has no syntax errors. It has no syntax errors. Therefore, it compiles.
  \item (b) If the network is down, then the site is unavailable. The site is
        available. Therefore, the network is not down.
\end{itemize}
\end{CoreProblem}
\begin{CoreProblem}{B.5}{Direct proof}
Let $A,B,C$ be sets. Give a direct element proof that if $A\subseteq B$, then
$A\cap C\subseteq B\cap C$.
\end{CoreProblem}
\begin{CoreProblem}{C.1}{Contraposition or contradiction}
Let $m,n\in\mathbb Z$. Prove by contraposition that if $mn$ is odd, then both
$m$ and $n$ are odd. Your proof must state the contrapositive explicitly and
handle its disjunction by cases.
\end{CoreProblem}
\columnbreak
\begin{CoreProblem}{D.3}{Cases, existence, or counterexample}
Give a constructive existence proof that there are four consecutive composite
positive integers. Verify explicitly that every integer in your example is
composite.
\end{CoreProblem}
\begin{CoreProblem}{E.2}{Strategy choice and repair}
The following proof is intended to show that the sum of two even integers is even. Identify the gap and repair it.

``Let $a$ and $b$ be even. Then $a=2k$ and $b=2k$ for some integer $k$. Hence $a+b=4k$, which is even.''
\end{CoreProblem}
\end{multicols}
\CoreWorkArea
\FormBreak
\CoreHeader{Proofs}{P36}{A--E}{\FivePartScoring Use complete sentences, name the requested method or rule, and justify every inference.}
\begin{multicols}{2}
\begin{CoreProblem}{A.3}{Inference and proof structure}
Over one common domain, identify the error in the proposed inference. Then give
a two-element domain and interpretations of $P$ and $Q$ that make both premises
true and the conclusion false.
\[
\exists x\,P(x),\qquad \exists x\,Q(x)
\quad\therefore\quad
\exists x\,(P(x)\land Q(x)).
\]
\end{CoreProblem}
\begin{CoreProblem}{B.2}{Direct proof}
Give a direct proof that the sum of two rational numbers is rational.
\end{CoreProblem}
\begin{CoreProblem}{C.1}{Contraposition or contradiction}
Let $m,n\in\mathbb Z$. Prove by contraposition that if $mn$ is odd, then both
$m$ and $n$ are odd. Your proof must state the contrapositive explicitly and
handle its disjunction by cases.
\end{CoreProblem}
\columnbreak
\begin{CoreProblem}{D.2}{Cases, existence, or counterexample}
Prove by cases that $|x|\ge 0$ for every real number $x$, with equality exactly when $x=0$.
\end{CoreProblem}
\begin{CoreProblem}{E.4}{Strategy choice and repair}
The following argument is circular. Explain the circularity, then replace it with a valid proof.

``To prove that $n^2$ is odd whenever $n$ is odd, assume $n^2$ is odd. Thus
$n^2=2k+1$ for some integer $k$, so $n^2$ is odd.''
\end{CoreProblem}
\end{multicols}
\CoreWorkArea
\FormBreak
\CoreHeader{Proofs}{P37}{A--E}{\FivePartScoring Use complete sentences, name the requested method or rule, and justify every inference.}
\begin{multicols}{2}
\begin{CoreProblem}{A.5}{Inference and proof structure}
Over a nonempty domain, use the premises
$\forall x\,(P(x)\to Q(x))$ and $\forall x\,P(x)$ to derive
$\forall x\,Q(x)$. Name the rule used at each step.
\end{CoreProblem}
\begin{CoreProblem}{B.2}{Direct proof}
Give a direct proof that the sum of two rational numbers is rational.
\end{CoreProblem}
\begin{CoreProblem}{C.2}{Contraposition or contradiction}
Prove by contraposition that if $3n+2$ is odd, then the integer $n$ is odd.
\end{CoreProblem}
\columnbreak
\begin{CoreProblem}{D.5}{Cases, existence, or counterexample}
Give a constructive existence proof that there are two irrational real numbers
whose product is rational.
\end{CoreProblem}
\begin{CoreProblem}{E.4}{Strategy choice and repair}
The following argument is circular. Explain the circularity, then replace it with a valid proof.

``To prove that $n^2$ is odd whenever $n$ is odd, assume $n^2$ is odd. Thus
$n^2=2k+1$ for some integer $k$, so $n^2$ is odd.''
\end{CoreProblem}
\end{multicols}
\CoreWorkArea
\FormBreak
\CoreHeader{Proofs}{P38}{A--E}{\FivePartScoring Use complete sentences, name the requested method or rule, and justify every inference.}
\begin{multicols}{2}
\begin{CoreProblem}{A.1}{Inference and proof structure}
Suppose that if the server is overloaded, requests are delayed; if requests are
delayed, the monitoring system sends an alert; and the server is overloaded.
Derive that the monitoring system sends an alert, naming the rule used at each
step.
\end{CoreProblem}
\begin{CoreProblem}{B.5}{Direct proof}
Let $A,B,C$ be sets. Give a direct element proof that if $A\subseteq B$, then
$A\cap C\subseteq B\cap C$.
\end{CoreProblem}
\begin{CoreProblem}{C.2}{Contraposition or contradiction}
Prove by contraposition that if $3n+2$ is odd, then the integer $n$ is odd.
\end{CoreProblem}
\columnbreak
\begin{CoreProblem}{D.5}{Cases, existence, or counterexample}
Give a constructive existence proof that there are two irrational real numbers
whose product is rational.
\end{CoreProblem}
\begin{CoreProblem}{E.2}{Strategy choice and repair}
The following proof is intended to show that the sum of two even integers is even. Identify the gap and repair it.

``Let $a$ and $b$ be even. Then $a=2k$ and $b=2k$ for some integer $k$. Hence $a+b=4k$, which is even.''
\end{CoreProblem}
\end{multicols}
\CoreWorkArea
\FormBreak
\CoreHeader{Proofs}{P39}{A--E}{\FivePartScoring Use complete sentences, name the requested method or rule, and justify every inference.}
\begin{multicols}{2}
\begin{CoreProblem}{A.2}{Inference and proof structure}
Determine whether each argument is valid. Name the valid inference rule or the invalid form.
\begin{itemize}
  \item (a) If the program compiles, then it has no syntax errors. It has no syntax errors. Therefore, it compiles.
  \item (b) If the network is down, then the site is unavailable. The site is
        available. Therefore, the network is not down.
\end{itemize}
\end{CoreProblem}
\begin{CoreProblem}{B.1}{Direct proof}
Give a direct proof that if an integer $n$ is even, then $n^2$ is even.
\end{CoreProblem}
\begin{CoreProblem}{C.4}{Contraposition or contradiction}
Let $r$ be rational and $x$ irrational. Prove by contradiction that $r+x$ is irrational.
\end{CoreProblem}
\columnbreak
\begin{CoreProblem}{D.5}{Cases, existence, or counterexample}
Give a constructive existence proof that there are two irrational real numbers
whose product is rational.
\end{CoreProblem}
\begin{CoreProblem}{E.2}{Strategy choice and repair}
The following proof is intended to show that the sum of two even integers is even. Identify the gap and repair it.

``Let $a$ and $b$ be even. Then $a=2k$ and $b=2k$ for some integer $k$. Hence $a+b=4k$, which is even.''
\end{CoreProblem}
\end{multicols}
\CoreWorkArea
\FormBreak
\CoreHeader{Proofs}{P40}{A--E}{\FivePartScoring Use complete sentences, name the requested method or rule, and justify every inference.}
\begin{multicols}{2}
\begin{CoreProblem}{A.4}{Inference and proof structure}
From $p\to q$, $q\to r$, and $\neg r$, derive $\neg p$. Name the rule used at
each step.
\end{CoreProblem}
\begin{CoreProblem}{B.3}{Direct proof}
Give a direct proof that the sum of two odd integers is even.
\end{CoreProblem}
\begin{CoreProblem}{C.1}{Contraposition or contradiction}
Let $m,n\in\mathbb Z$. Prove by contraposition that if $mn$ is odd, then both
$m$ and $n$ are odd. Your proof must state the contrapositive explicitly and
handle its disjunction by cases.
\end{CoreProblem}
\columnbreak
\begin{CoreProblem}{D.3}{Cases, existence, or counterexample}
Give a constructive existence proof that there are four consecutive composite
positive integers. Verify explicitly that every integer in your example is
composite.
\end{CoreProblem}
\begin{CoreProblem}{E.3}{Strategy choice and repair}
The following proof is intended to show that the product of a rational number
and an irrational number is irrational. Identify why the argument fails when
$r=0$. Add the weakest necessary restriction on $r$, state the corrected
theorem, and prove it.

``Let $r$ be rational and $x$ irrational. If $rx$ were rational, then $x=(rx)/r$ would be rational, a contradiction.''
\end{CoreProblem}
\end{multicols}
\CoreWorkArea
\FormBreak
\CoreHeader{Proofs}{P41}{A--E}{\FivePartScoring Use complete sentences, name the requested method or rule, and justify every inference.}
\begin{multicols}{2}
\begin{CoreProblem}{A.5}{Inference and proof structure}
Over a nonempty domain, use the premises
$\forall x\,(P(x)\to Q(x))$ and $\forall x\,P(x)$ to derive
$\forall x\,Q(x)$. Name the rule used at each step.
\end{CoreProblem}
\begin{CoreProblem}{B.2}{Direct proof}
Give a direct proof that the sum of two rational numbers is rational.
\end{CoreProblem}
\begin{CoreProblem}{C.4}{Contraposition or contradiction}
Let $r$ be rational and $x$ irrational. Prove by contradiction that $r+x$ is irrational.
\end{CoreProblem}
\columnbreak
\begin{CoreProblem}{D.3}{Cases, existence, or counterexample}
Give a constructive existence proof that there are four consecutive composite
positive integers. Verify explicitly that every integer in your example is
composite.
\end{CoreProblem}
\begin{CoreProblem}{E.1}{Strategy choice and repair}
For each statement, choose a promising first strategy---direct proof,
contraposition, contradiction, cases, constructive existence, or
counterexample---and state the first substantive step. Do not write the full
proofs.
\begin{itemize}
  \item (a) If $n^2$ is even, then the integer $n$ is even.
  \item (b) There is no smallest positive real number.
\end{itemize}
\end{CoreProblem}
\end{multicols}
\CoreWorkArea
\FormBreak
\CoreHeader{Proofs}{P42}{A--E}{\FivePartScoring Use complete sentences, name the requested method or rule, and justify every inference.}
\begin{multicols}{2}
\begin{CoreProblem}{A.3}{Inference and proof structure}
Over one common domain, identify the error in the proposed inference. Then give
a two-element domain and interpretations of $P$ and $Q$ that make both premises
true and the conclusion false.
\[
\exists x\,P(x),\qquad \exists x\,Q(x)
\quad\therefore\quad
\exists x\,(P(x)\land Q(x)).
\]
\end{CoreProblem}
\begin{CoreProblem}{B.3}{Direct proof}
Give a direct proof that the sum of two odd integers is even.
\end{CoreProblem}
\begin{CoreProblem}{C.5}{Contraposition or contradiction}
Prove by contraposition that if $n^2$ is even, then the integer $n$ is even.
\end{CoreProblem}
\columnbreak
\begin{CoreProblem}{D.5}{Cases, existence, or counterexample}
Give a constructive existence proof that there are two irrational real numbers
whose product is rational.
\end{CoreProblem}
\begin{CoreProblem}{E.2}{Strategy choice and repair}
The following proof is intended to show that the sum of two even integers is even. Identify the gap and repair it.

``Let $a$ and $b$ be even. Then $a=2k$ and $b=2k$ for some integer $k$. Hence $a+b=4k$, which is even.''
\end{CoreProblem}
\end{multicols}
\CoreWorkArea
\FormBreak
\CoreHeader{Proofs}{P43}{A--E}{\FivePartScoring Use complete sentences, name the requested method or rule, and justify every inference.}
\begin{multicols}{2}
\begin{CoreProblem}{A.3}{Inference and proof structure}
Over one common domain, identify the error in the proposed inference. Then give
a two-element domain and interpretations of $P$ and $Q$ that make both premises
true and the conclusion false.
\[
\exists x\,P(x),\qquad \exists x\,Q(x)
\quad\therefore\quad
\exists x\,(P(x)\land Q(x)).
\]
\end{CoreProblem}
\begin{CoreProblem}{B.2}{Direct proof}
Give a direct proof that the sum of two rational numbers is rational.
\end{CoreProblem}
\begin{CoreProblem}{C.2}{Contraposition or contradiction}
Prove by contraposition that if $3n+2$ is odd, then the integer $n$ is odd.
\end{CoreProblem}
\columnbreak
\begin{CoreProblem}{D.5}{Cases, existence, or counterexample}
Give a constructive existence proof that there are two irrational real numbers
whose product is rational.
\end{CoreProblem}
\begin{CoreProblem}{E.1}{Strategy choice and repair}
For each statement, choose a promising first strategy---direct proof,
contraposition, contradiction, cases, constructive existence, or
counterexample---and state the first substantive step. Do not write the full
proofs.
\begin{itemize}
  \item (a) If $n^2$ is even, then the integer $n$ is even.
  \item (b) There is no smallest positive real number.
\end{itemize}
\end{CoreProblem}
\end{multicols}
\CoreWorkArea
\FormBreak
\CoreHeader{Proofs}{P44}{A--E}{\FivePartScoring Use complete sentences, name the requested method or rule, and justify every inference.}
\begin{multicols}{2}
\begin{CoreProblem}{A.3}{Inference and proof structure}
Over one common domain, identify the error in the proposed inference. Then give
a two-element domain and interpretations of $P$ and $Q$ that make both premises
true and the conclusion false.
\[
\exists x\,P(x),\qquad \exists x\,Q(x)
\quad\therefore\quad
\exists x\,(P(x)\land Q(x)).
\]
\end{CoreProblem}
\begin{CoreProblem}{B.5}{Direct proof}
Let $A,B,C$ be sets. Give a direct element proof that if $A\subseteq B$, then
$A\cap C\subseteq B\cap C$.
\end{CoreProblem}
\begin{CoreProblem}{C.4}{Contraposition or contradiction}
Let $r$ be rational and $x$ irrational. Prove by contradiction that $r+x$ is irrational.
\end{CoreProblem}
\columnbreak
\begin{CoreProblem}{D.5}{Cases, existence, or counterexample}
Give a constructive existence proof that there are two irrational real numbers
whose product is rational.
\end{CoreProblem}
\begin{CoreProblem}{E.1}{Strategy choice and repair}
For each statement, choose a promising first strategy---direct proof,
contraposition, contradiction, cases, constructive existence, or
counterexample---and state the first substantive step. Do not write the full
proofs.
\begin{itemize}
  \item (a) If $n^2$ is even, then the integer $n$ is even.
  \item (b) There is no smallest positive real number.
\end{itemize}
\end{CoreProblem}
\end{multicols}
\CoreWorkArea
\FormBreak
\CoreHeader{Proofs}{P45}{A--E}{\FivePartScoring Use complete sentences, name the requested method or rule, and justify every inference.}
\begin{multicols}{2}
\begin{CoreProblem}{A.1}{Inference and proof structure}
Suppose that if the server is overloaded, requests are delayed; if requests are
delayed, the monitoring system sends an alert; and the server is overloaded.
Derive that the monitoring system sends an alert, naming the rule used at each
step.
\end{CoreProblem}
\begin{CoreProblem}{B.1}{Direct proof}
Give a direct proof that if an integer $n$ is even, then $n^2$ is even.
\end{CoreProblem}
\begin{CoreProblem}{C.2}{Contraposition or contradiction}
Prove by contraposition that if $3n+2$ is odd, then the integer $n$ is odd.
\end{CoreProblem}
\columnbreak
\begin{CoreProblem}{D.4}{Cases, existence, or counterexample}
Disprove each universal statement by giving a counterexample.
\begin{itemize}
  \item (a) The sum of two irrational numbers is irrational.
  \item (b) $n^2+n+41$ is prime for every nonnegative integer $n$.
\end{itemize}
\end{CoreProblem}
\begin{CoreProblem}{E.5}{Strategy choice and repair}
The following proof is intended to show that $A\subseteq B$ whenever
$A\cap B=A$. Identify the gap and repair it.

``Let $x\in A\cap B$. Then $x\in B$. Therefore $A\subseteq B$.''
\end{CoreProblem}
\end{multicols}
\CoreWorkArea
\FormBreak
\CoreHeader{Proofs}{P46}{A--E}{\FivePartScoring Use complete sentences, name the requested method or rule, and justify every inference.}
\begin{multicols}{2}
\begin{CoreProblem}{A.5}{Inference and proof structure}
Over a nonempty domain, use the premises
$\forall x\,(P(x)\to Q(x))$ and $\forall x\,P(x)$ to derive
$\forall x\,Q(x)$. Name the rule used at each step.
\end{CoreProblem}
\begin{CoreProblem}{B.1}{Direct proof}
Give a direct proof that if an integer $n$ is even, then $n^2$ is even.
\end{CoreProblem}
\begin{CoreProblem}{C.2}{Contraposition or contradiction}
Prove by contraposition that if $3n+2$ is odd, then the integer $n$ is odd.
\end{CoreProblem}
\columnbreak
\begin{CoreProblem}{D.4}{Cases, existence, or counterexample}
Disprove each universal statement by giving a counterexample.
\begin{itemize}
  \item (a) The sum of two irrational numbers is irrational.
  \item (b) $n^2+n+41$ is prime for every nonnegative integer $n$.
\end{itemize}
\end{CoreProblem}
\begin{CoreProblem}{E.5}{Strategy choice and repair}
The following proof is intended to show that $A\subseteq B$ whenever
$A\cap B=A$. Identify the gap and repair it.

``Let $x\in A\cap B$. Then $x\in B$. Therefore $A\subseteq B$.''
\end{CoreProblem}
\end{multicols}
\CoreWorkArea
\FormBreak
\CoreHeader{Proofs}{P47}{A--E}{\FivePartScoring Use complete sentences, name the requested method or rule, and justify every inference.}
\begin{multicols}{2}
\begin{CoreProblem}{A.5}{Inference and proof structure}
Over a nonempty domain, use the premises
$\forall x\,(P(x)\to Q(x))$ and $\forall x\,P(x)$ to derive
$\forall x\,Q(x)$. Name the rule used at each step.
\end{CoreProblem}
\begin{CoreProblem}{B.2}{Direct proof}
Give a direct proof that the sum of two rational numbers is rational.
\end{CoreProblem}
\begin{CoreProblem}{C.4}{Contraposition or contradiction}
Let $r$ be rational and $x$ irrational. Prove by contradiction that $r+x$ is irrational.
\end{CoreProblem}
\columnbreak
\begin{CoreProblem}{D.3}{Cases, existence, or counterexample}
Give a constructive existence proof that there are four consecutive composite
positive integers. Verify explicitly that every integer in your example is
composite.
\end{CoreProblem}
\begin{CoreProblem}{E.3}{Strategy choice and repair}
The following proof is intended to show that the product of a rational number
and an irrational number is irrational. Identify why the argument fails when
$r=0$. Add the weakest necessary restriction on $r$, state the corrected
theorem, and prove it.

``Let $r$ be rational and $x$ irrational. If $rx$ were rational, then $x=(rx)/r$ would be rational, a contradiction.''
\end{CoreProblem}
\end{multicols}
\CoreWorkArea
\FormBreak
\CoreHeader{Proofs}{P48}{A--E}{\FivePartScoring Use complete sentences, name the requested method or rule, and justify every inference.}
\begin{multicols}{2}
\begin{CoreProblem}{A.5}{Inference and proof structure}
Over a nonempty domain, use the premises
$\forall x\,(P(x)\to Q(x))$ and $\forall x\,P(x)$ to derive
$\forall x\,Q(x)$. Name the rule used at each step.
\end{CoreProblem}
\begin{CoreProblem}{B.2}{Direct proof}
Give a direct proof that the sum of two rational numbers is rational.
\end{CoreProblem}
\begin{CoreProblem}{C.1}{Contraposition or contradiction}
Let $m,n\in\mathbb Z$. Prove by contraposition that if $mn$ is odd, then both
$m$ and $n$ are odd. Your proof must state the contrapositive explicitly and
handle its disjunction by cases.
\end{CoreProblem}
\columnbreak
\begin{CoreProblem}{D.1}{Cases, existence, or counterexample}
Prove by cases that $n(n+1)$ is even for every integer $n$.
\end{CoreProblem}
\begin{CoreProblem}{E.4}{Strategy choice and repair}
The following argument is circular. Explain the circularity, then replace it with a valid proof.

``To prove that $n^2$ is odd whenever $n$ is odd, assume $n^2$ is odd. Thus
$n^2=2k+1$ for some integer $k$, so $n^2$ is odd.''
\end{CoreProblem}
\end{multicols}
\CoreWorkArea
\FormBreak
\CoreHeader{Proofs}{P49}{A--E}{\FivePartScoring Use complete sentences, name the requested method or rule, and justify every inference.}
\begin{multicols}{2}
\begin{CoreProblem}{A.5}{Inference and proof structure}
Over a nonempty domain, use the premises
$\forall x\,(P(x)\to Q(x))$ and $\forall x\,P(x)$ to derive
$\forall x\,Q(x)$. Name the rule used at each step.
\end{CoreProblem}
\begin{CoreProblem}{B.5}{Direct proof}
Let $A,B,C$ be sets. Give a direct element proof that if $A\subseteq B$, then
$A\cap C\subseteq B\cap C$.
\end{CoreProblem}
\begin{CoreProblem}{C.1}{Contraposition or contradiction}
Let $m,n\in\mathbb Z$. Prove by contraposition that if $mn$ is odd, then both
$m$ and $n$ are odd. Your proof must state the contrapositive explicitly and
handle its disjunction by cases.
\end{CoreProblem}
\columnbreak
\begin{CoreProblem}{D.3}{Cases, existence, or counterexample}
Give a constructive existence proof that there are four consecutive composite
positive integers. Verify explicitly that every integer in your example is
composite.
\end{CoreProblem}
\begin{CoreProblem}{E.5}{Strategy choice and repair}
The following proof is intended to show that $A\subseteq B$ whenever
$A\cap B=A$. Identify the gap and repair it.

``Let $x\in A\cap B$. Then $x\in B$. Therefore $A\subseteq B$.''
\end{CoreProblem}
\end{multicols}
\CoreWorkArea
\FormBreak
\CoreHeader{Proofs}{P50}{A--E}{\FivePartScoring Use complete sentences, name the requested method or rule, and justify every inference.}
\begin{multicols}{2}
\begin{CoreProblem}{A.5}{Inference and proof structure}
Over a nonempty domain, use the premises
$\forall x\,(P(x)\to Q(x))$ and $\forall x\,P(x)$ to derive
$\forall x\,Q(x)$. Name the rule used at each step.
\end{CoreProblem}
\begin{CoreProblem}{B.2}{Direct proof}
Give a direct proof that the sum of two rational numbers is rational.
\end{CoreProblem}
\begin{CoreProblem}{C.4}{Contraposition or contradiction}
Let $r$ be rational and $x$ irrational. Prove by contradiction that $r+x$ is irrational.
\end{CoreProblem}
\columnbreak
\begin{CoreProblem}{D.2}{Cases, existence, or counterexample}
Prove by cases that $|x|\ge 0$ for every real number $x$, with equality exactly when $x=0$.
\end{CoreProblem}
\begin{CoreProblem}{E.2}{Strategy choice and repair}
The following proof is intended to show that the sum of two even integers is even. Identify the gap and repair it.

``Let $a$ and $b$ be even. Then $a=2k$ and $b=2k$ for some integer $k$. Hence $a+b=4k$, which is even.''
\end{CoreProblem}
\end{multicols}
\CoreWorkArea
\FormBreak
\CoreHeader{Proofs}{P51}{A--E}{\FivePartScoring Use complete sentences, name the requested method or rule, and justify every inference.}
\begin{multicols}{2}
\begin{CoreProblem}{A.5}{Inference and proof structure}
Over a nonempty domain, use the premises
$\forall x\,(P(x)\to Q(x))$ and $\forall x\,P(x)$ to derive
$\forall x\,Q(x)$. Name the rule used at each step.
\end{CoreProblem}
\begin{CoreProblem}{B.5}{Direct proof}
Let $A,B,C$ be sets. Give a direct element proof that if $A\subseteq B$, then
$A\cap C\subseteq B\cap C$.
\end{CoreProblem}
\begin{CoreProblem}{C.4}{Contraposition or contradiction}
Let $r$ be rational and $x$ irrational. Prove by contradiction that $r+x$ is irrational.
\end{CoreProblem}
\columnbreak
\begin{CoreProblem}{D.1}{Cases, existence, or counterexample}
Prove by cases that $n(n+1)$ is even for every integer $n$.
\end{CoreProblem}
\begin{CoreProblem}{E.1}{Strategy choice and repair}
For each statement, choose a promising first strategy---direct proof,
contraposition, contradiction, cases, constructive existence, or
counterexample---and state the first substantive step. Do not write the full
proofs.
\begin{itemize}
  \item (a) If $n^2$ is even, then the integer $n$ is even.
  \item (b) There is no smallest positive real number.
\end{itemize}
\end{CoreProblem}
\end{multicols}
\CoreWorkArea
\FormBreak
\CoreHeader{Proofs}{P52}{A--E}{\FivePartScoring Use complete sentences, name the requested method or rule, and justify every inference.}
\begin{multicols}{2}
\begin{CoreProblem}{A.2}{Inference and proof structure}
Determine whether each argument is valid. Name the valid inference rule or the invalid form.
\begin{itemize}
  \item (a) If the program compiles, then it has no syntax errors. It has no syntax errors. Therefore, it compiles.
  \item (b) If the network is down, then the site is unavailable. The site is
        available. Therefore, the network is not down.
\end{itemize}
\end{CoreProblem}
\begin{CoreProblem}{B.3}{Direct proof}
Give a direct proof that the sum of two odd integers is even.
\end{CoreProblem}
\begin{CoreProblem}{C.5}{Contraposition or contradiction}
Prove by contraposition that if $n^2$ is even, then the integer $n$ is even.
\end{CoreProblem}
\columnbreak
\begin{CoreProblem}{D.3}{Cases, existence, or counterexample}
Give a constructive existence proof that there are four consecutive composite
positive integers. Verify explicitly that every integer in your example is
composite.
\end{CoreProblem}
\begin{CoreProblem}{E.2}{Strategy choice and repair}
The following proof is intended to show that the sum of two even integers is even. Identify the gap and repair it.

``Let $a$ and $b$ be even. Then $a=2k$ and $b=2k$ for some integer $k$. Hence $a+b=4k$, which is even.''
\end{CoreProblem}
\end{multicols}
\CoreWorkArea
\FormBreak
\CoreHeader{Proofs}{P53}{A--E}{\FivePartScoring Use complete sentences, name the requested method or rule, and justify every inference.}
\begin{multicols}{2}
\begin{CoreProblem}{A.1}{Inference and proof structure}
Suppose that if the server is overloaded, requests are delayed; if requests are
delayed, the monitoring system sends an alert; and the server is overloaded.
Derive that the monitoring system sends an alert, naming the rule used at each
step.
\end{CoreProblem}
\begin{CoreProblem}{B.4}{Direct proof}
Give a direct proof that the difference of two even integers is even.
\end{CoreProblem}
\begin{CoreProblem}{C.5}{Contraposition or contradiction}
Prove by contraposition that if $n^2$ is even, then the integer $n$ is even.
\end{CoreProblem}
\columnbreak
\begin{CoreProblem}{D.2}{Cases, existence, or counterexample}
Prove by cases that $|x|\ge 0$ for every real number $x$, with equality exactly when $x=0$.
\end{CoreProblem}
\begin{CoreProblem}{E.3}{Strategy choice and repair}
The following proof is intended to show that the product of a rational number
and an irrational number is irrational. Identify why the argument fails when
$r=0$. Add the weakest necessary restriction on $r$, state the corrected
theorem, and prove it.

``Let $r$ be rational and $x$ irrational. If $rx$ were rational, then $x=(rx)/r$ would be rational, a contradiction.''
\end{CoreProblem}
\end{multicols}
\CoreWorkArea
\FormBreak
\CoreHeader{Proofs}{P54}{A--E}{\FivePartScoring Use complete sentences, name the requested method or rule, and justify every inference.}
\begin{multicols}{2}
\begin{CoreProblem}{A.2}{Inference and proof structure}
Determine whether each argument is valid. Name the valid inference rule or the invalid form.
\begin{itemize}
  \item (a) If the program compiles, then it has no syntax errors. It has no syntax errors. Therefore, it compiles.
  \item (b) If the network is down, then the site is unavailable. The site is
        available. Therefore, the network is not down.
\end{itemize}
\end{CoreProblem}
\begin{CoreProblem}{B.5}{Direct proof}
Let $A,B,C$ be sets. Give a direct element proof that if $A\subseteq B$, then
$A\cap C\subseteq B\cap C$.
\end{CoreProblem}
\begin{CoreProblem}{C.4}{Contraposition or contradiction}
Let $r$ be rational and $x$ irrational. Prove by contradiction that $r+x$ is irrational.
\end{CoreProblem}
\columnbreak
\begin{CoreProblem}{D.1}{Cases, existence, or counterexample}
Prove by cases that $n(n+1)$ is even for every integer $n$.
\end{CoreProblem}
\begin{CoreProblem}{E.1}{Strategy choice and repair}
For each statement, choose a promising first strategy---direct proof,
contraposition, contradiction, cases, constructive existence, or
counterexample---and state the first substantive step. Do not write the full
proofs.
\begin{itemize}
  \item (a) If $n^2$ is even, then the integer $n$ is even.
  \item (b) There is no smallest positive real number.
\end{itemize}
\end{CoreProblem}
\end{multicols}
\CoreWorkArea
\FormBreak
\CoreHeader{Proofs}{P55}{A--E}{\FivePartScoring Use complete sentences, name the requested method or rule, and justify every inference.}
\begin{multicols}{2}
\begin{CoreProblem}{A.1}{Inference and proof structure}
Suppose that if the server is overloaded, requests are delayed; if requests are
delayed, the monitoring system sends an alert; and the server is overloaded.
Derive that the monitoring system sends an alert, naming the rule used at each
step.
\end{CoreProblem}
\begin{CoreProblem}{B.2}{Direct proof}
Give a direct proof that the sum of two rational numbers is rational.
\end{CoreProblem}
\begin{CoreProblem}{C.5}{Contraposition or contradiction}
Prove by contraposition that if $n^2$ is even, then the integer $n$ is even.
\end{CoreProblem}
\columnbreak
\begin{CoreProblem}{D.3}{Cases, existence, or counterexample}
Give a constructive existence proof that there are four consecutive composite
positive integers. Verify explicitly that every integer in your example is
composite.
\end{CoreProblem}
\begin{CoreProblem}{E.3}{Strategy choice and repair}
The following proof is intended to show that the product of a rational number
and an irrational number is irrational. Identify why the argument fails when
$r=0$. Add the weakest necessary restriction on $r$, state the corrected
theorem, and prove it.

``Let $r$ be rational and $x$ irrational. If $rx$ were rational, then $x=(rx)/r$ would be rational, a contradiction.''
\end{CoreProblem}
\end{multicols}
\CoreWorkArea
\FormBreak
\CoreHeader{Proofs}{P56}{A--E}{\FivePartScoring Use complete sentences, name the requested method or rule, and justify every inference.}
\begin{multicols}{2}
\begin{CoreProblem}{A.4}{Inference and proof structure}
From $p\to q$, $q\to r$, and $\neg r$, derive $\neg p$. Name the rule used at
each step.
\end{CoreProblem}
\begin{CoreProblem}{B.2}{Direct proof}
Give a direct proof that the sum of two rational numbers is rational.
\end{CoreProblem}
\begin{CoreProblem}{C.1}{Contraposition or contradiction}
Let $m,n\in\mathbb Z$. Prove by contraposition that if $mn$ is odd, then both
$m$ and $n$ are odd. Your proof must state the contrapositive explicitly and
handle its disjunction by cases.
\end{CoreProblem}
\columnbreak
\begin{CoreProblem}{D.2}{Cases, existence, or counterexample}
Prove by cases that $|x|\ge 0$ for every real number $x$, with equality exactly when $x=0$.
\end{CoreProblem}
\begin{CoreProblem}{E.4}{Strategy choice and repair}
The following argument is circular. Explain the circularity, then replace it with a valid proof.

``To prove that $n^2$ is odd whenever $n$ is odd, assume $n^2$ is odd. Thus
$n^2=2k+1$ for some integer $k$, so $n^2$ is odd.''
\end{CoreProblem}
\end{multicols}
\CoreWorkArea
\FormBreak
\CoreHeader{Proofs}{P57}{A--E}{\FivePartScoring Use complete sentences, name the requested method or rule, and justify every inference.}
\begin{multicols}{2}
\begin{CoreProblem}{A.3}{Inference and proof structure}
Over one common domain, identify the error in the proposed inference. Then give
a two-element domain and interpretations of $P$ and $Q$ that make both premises
true and the conclusion false.
\[
\exists x\,P(x),\qquad \exists x\,Q(x)
\quad\therefore\quad
\exists x\,(P(x)\land Q(x)).
\]
\end{CoreProblem}
\begin{CoreProblem}{B.5}{Direct proof}
Let $A,B,C$ be sets. Give a direct element proof that if $A\subseteq B$, then
$A\cap C\subseteq B\cap C$.
\end{CoreProblem}
\begin{CoreProblem}{C.5}{Contraposition or contradiction}
Prove by contraposition that if $n^2$ is even, then the integer $n$ is even.
\end{CoreProblem}
\columnbreak
\begin{CoreProblem}{D.4}{Cases, existence, or counterexample}
Disprove each universal statement by giving a counterexample.
\begin{itemize}
  \item (a) The sum of two irrational numbers is irrational.
  \item (b) $n^2+n+41$ is prime for every nonnegative integer $n$.
\end{itemize}
\end{CoreProblem}
\begin{CoreProblem}{E.5}{Strategy choice and repair}
The following proof is intended to show that $A\subseteq B$ whenever
$A\cap B=A$. Identify the gap and repair it.

``Let $x\in A\cap B$. Then $x\in B$. Therefore $A\subseteq B$.''
\end{CoreProblem}
\end{multicols}
\CoreWorkArea
\FormBreak
\CoreHeader{Proofs}{P58}{A--E}{\FivePartScoring Use complete sentences, name the requested method or rule, and justify every inference.}
\begin{multicols}{2}
\begin{CoreProblem}{A.1}{Inference and proof structure}
Suppose that if the server is overloaded, requests are delayed; if requests are
delayed, the monitoring system sends an alert; and the server is overloaded.
Derive that the monitoring system sends an alert, naming the rule used at each
step.
\end{CoreProblem}
\begin{CoreProblem}{B.4}{Direct proof}
Give a direct proof that the difference of two even integers is even.
\end{CoreProblem}
\begin{CoreProblem}{C.2}{Contraposition or contradiction}
Prove by contraposition that if $3n+2$ is odd, then the integer $n$ is odd.
\end{CoreProblem}
\columnbreak
\begin{CoreProblem}{D.1}{Cases, existence, or counterexample}
Prove by cases that $n(n+1)$ is even for every integer $n$.
\end{CoreProblem}
\begin{CoreProblem}{E.1}{Strategy choice and repair}
For each statement, choose a promising first strategy---direct proof,
contraposition, contradiction, cases, constructive existence, or
counterexample---and state the first substantive step. Do not write the full
proofs.
\begin{itemize}
  \item (a) If $n^2$ is even, then the integer $n$ is even.
  \item (b) There is no smallest positive real number.
\end{itemize}
\end{CoreProblem}
\end{multicols}
\CoreWorkArea
\FormBreak
\CoreHeader{Proofs}{P59}{A--E}{\FivePartScoring Use complete sentences, name the requested method or rule, and justify every inference.}
\begin{multicols}{2}
\begin{CoreProblem}{A.3}{Inference and proof structure}
Over one common domain, identify the error in the proposed inference. Then give
a two-element domain and interpretations of $P$ and $Q$ that make both premises
true and the conclusion false.
\[
\exists x\,P(x),\qquad \exists x\,Q(x)
\quad\therefore\quad
\exists x\,(P(x)\land Q(x)).
\]
\end{CoreProblem}
\begin{CoreProblem}{B.5}{Direct proof}
Let $A,B,C$ be sets. Give a direct element proof that if $A\subseteq B$, then
$A\cap C\subseteq B\cap C$.
\end{CoreProblem}
\begin{CoreProblem}{C.1}{Contraposition or contradiction}
Let $m,n\in\mathbb Z$. Prove by contraposition that if $mn$ is odd, then both
$m$ and $n$ are odd. Your proof must state the contrapositive explicitly and
handle its disjunction by cases.
\end{CoreProblem}
\columnbreak
\begin{CoreProblem}{D.3}{Cases, existence, or counterexample}
Give a constructive existence proof that there are four consecutive composite
positive integers. Verify explicitly that every integer in your example is
composite.
\end{CoreProblem}
\begin{CoreProblem}{E.1}{Strategy choice and repair}
For each statement, choose a promising first strategy---direct proof,
contraposition, contradiction, cases, constructive existence, or
counterexample---and state the first substantive step. Do not write the full
proofs.
\begin{itemize}
  \item (a) If $n^2$ is even, then the integer $n$ is even.
  \item (b) There is no smallest positive real number.
\end{itemize}
\end{CoreProblem}
\end{multicols}
\CoreWorkArea
\FormBreak
\CoreHeader{Proofs}{P60}{A--E}{\FivePartScoring Use complete sentences, name the requested method or rule, and justify every inference.}
\begin{multicols}{2}
\begin{CoreProblem}{A.4}{Inference and proof structure}
From $p\to q$, $q\to r$, and $\neg r$, derive $\neg p$. Name the rule used at
each step.
\end{CoreProblem}
\begin{CoreProblem}{B.4}{Direct proof}
Give a direct proof that the difference of two even integers is even.
\end{CoreProblem}
\begin{CoreProblem}{C.3}{Contraposition or contradiction}
Prove by contradiction that there is no largest integer.
\end{CoreProblem}
\columnbreak
\begin{CoreProblem}{D.3}{Cases, existence, or counterexample}
Give a constructive existence proof that there are four consecutive composite
positive integers. Verify explicitly that every integer in your example is
composite.
\end{CoreProblem}
\begin{CoreProblem}{E.3}{Strategy choice and repair}
The following proof is intended to show that the product of a rational number
and an irrational number is irrational. Identify why the argument fails when
$r=0$. Add the weakest necessary restriction on $r$, state the corrected
theorem, and prove it.

``Let $r$ be rational and $x$ irrational. If $rx$ were rational, then $x=(rx)/r$ would be rational, a contradiction.''
\end{CoreProblem}
\end{multicols}
\CoreWorkArea
\end{document}
